\section{Coherence Diagnostics}\label{sec:tci}

The coherence-delay law constrains the representation term of the finite-memory
averaging systems studied here, but closed-loop performance also depends on
whether actions remain reproducible and whether adaptation remains effective.
The ISS template of \cref{sec:sufficiency} provides one way to score those
three ingredients. This section packages them into a diagnostic built from the
minimal decomposition $(\Pot,\Act,\Conv)$.

\begin{definition}[Component Coherence Scores]\label{def:component_scores}
Let $\IS = (\Pot,\Act,\Conv)$ be an information system with ISS-Lyapunov
functionals $V_P$, $V_A$, $V_\Phi$ as in \cref{sec:sufficiency}.
Define the three \emph{component coherence scores} at time $t$ as:
\begin{align}
  \sigma_P(t)
    &\coloneqq \frac{1}{1 + V_P(t)}
     = \frac{1}{1 + \tfrac{1}{2}W_2\!\left(\Pot(S,t),\,\Pot^*(S,t)\right)^2},
    \label{eq:sigmaP}\\[4pt]
  \sigma_A(t)
    &\coloneqq \exp\!\left(-H\!\left(\Act_t\mid\Pot(S,t)\right)\right),
    \label{eq:sigmaA}\\[4pt]
  \sigma_\Phi(t)
    &\coloneqq \frac{1}{1 + V_\Phi(t)}.
    \label{eq:sigmaPhi}
\end{align}
Each score takes values in $(0,1]$, equals $1$ when the corresponding
component is at its reference state, and approaches $0$ as the component degrades.
\end{definition}

\begin{definition}[Coherence Index]\label{def:tci}
The \emph{Coherence Index (CI)} of a system $\IS$ at time $t$, written formally
as $\Score(\IS,t)$ when the functional dependence must be shown explicitly, is
\[
  \Score(\IS, t)
  \;\coloneqq\;
  \min\!\left(\sigma_P(t),\;\sigma_A(t),\;\sigma_\Phi(t)\right)
  \;\in\;(0,1].
\]
\end{definition}

The $\min$ aggregation is the correct choice given the Borromean property
(\cref{thm:borromean}): system coherence is bounded by the weakest component,
and strength in two components does not compensate for failure in the third.

\begin{remark}[Bottleneck sensitivity and effort correction]
The use of $\min(\sigma_P, \sigma_A, \sigma_\Phi)$ is bottleneck-sensitive:
it diagnoses the weakest active subsystem, not overall system quality in the
round. In action-effective environments, high actuation capacity can keep
$\sigma_P$ artificially high by forcing the world toward a stale model.
When such feedback masking is possible, a conservative extension is the
\emph{effort-corrected score}
\[
  \ScoreE(\IS,t)
  \;\coloneqq\;
  \min\!\left(
    \sigma_P(t)e^{-\lambda E_t/E_0},\;\sigma_A(t),\;\sigma_\Phi(t)
  \right),
\]
where $E_t$ is any monotone proxy for corrective actuation burden,
$E_0 > 0$ is a reference scale, and $\lambda > 0$ controls sensitivity. It
discounts apparent improvements in model-world alignment when those gains come
from action-induced confounding rather than better state estimation. In the
synthetic particle tracker, $E_t$ is the realized coercive effort
$\alpha |a_t - s_t^{\mathrm{free}}|$; in deployed systems, action increments,
energy expenditure, or control-surface usage can play the same role.
For logged recommendation benchmarks, the preferred instantiation is the
per-user KL divergence between the current exposure-tag distribution and the
healthy random-exposure baseline distribution.
\end{remark}

\subsection{Estimable Score}

\Cref{def:component_scores} requires $\Pot^*(S,t)$, the true oracle
distribution of the environment, to compute $\sigma_P(t)$.
In any deployed system, $\Pot^*$ is unobservable by construction: a system
with complete knowledge of $\Pot^*$ would have nothing to estimate.
We therefore introduce an estimable version of $\sigma_P$ using an
empirical approximation of $\Pot^*$.

\begin{definition}[Empirical Potentiality Estimator]\label{def:potestimator}
Let $\{s_{t-n+1}, \ldots, s_t\}$ be the $n$ most recent environmental
signals observed by the system.
The \emph{empirical potentiality estimator} is the empirical measure:
\[
  \hat{\Pot}^*(S,t) \;\coloneqq\; \frac{1}{n} \sum_{k=t-n+1}^{t} \delta_{s_k},
\]
where $\delta_{s_k}$ is the Dirac measure at observation $s_k$.
\end{definition}

\begin{definition}[Estimable Score]\label{def:tci_estimable}
The \emph{debiased Sinkhorn divergence} between measures $\mu$, $\nu$ is
\[
  S_\varepsilon(\mu,\nu)
  \;\coloneqq\;
  T_\varepsilon(\mu,\nu)
  - \tfrac{1}{2}T_\varepsilon(\mu,\mu)
  - \tfrac{1}{2}T_\varepsilon(\nu,\nu),
\]
where $T_\varepsilon$ is the raw entropically regularized optimal transport cost.
Unlike $T_\varepsilon$, the debiased divergence satisfies $S_\varepsilon(\mu,\mu) = 0$,
restoring the metric property and enabling the clean exponent form of
\cref{prop:tci_estimation_error}, with ambient dimension, regularization, and
support geometry absorbed into the constants rather than the displayed rate.

The \emph{estimable component $\hat\sigma_P$} replaces $\Pot^*$ with
$\hat\Pot^*$ in \cref{eq:sigmaP} using $S_\varepsilon$:
\[
  \hat\sigma_P(t) \;\coloneqq\;
  \frac{1}{1 + \tfrac{1}{2}
    S_\varepsilon\!\left(\Pot(S,t),\,\hat{\Pot}^*(S,t)\right)^2}.
\]
The \emph{estimable score} is
\[
  \ScoreHat(\IS,t) \;\coloneqq\;
  \min\!\left(\hat\sigma_P(t),\;\sigma_A(t),\;\sigma_\Phi(t)\right).
\]
Operationally, the three components sit on an observability ladder: `exact`
when the action channel or Lyapunov functional is constructively known,
`identified` when logging is rich enough for a direct estimator, and `proxy`
when the benchmark supports only an operational surrogate. In the present
experiments, the Gaussian tracker is exact, the particle tracker is near the
identified/surrogate boundary, and KuaiRand is proxy-level.
\end{definition}

\begin{remark}[Null and alternative calibration for Sinkhorn scores]
Recent limit theorems for the empirical Sinkhorn divergence make an important
calibration split explicit \citep{GoldfeldEtAl2022}. Under the null $P = P^*$,
the debiased statistic has an $n^{-1}$ null fluctuation order, whereas under the
drifted alternative regime in which the present estimation bound is used the
statistical term follows the usual $n^{-1/2}$ scale. The practical implication
is methodological rather than asymptotic: healthy reference windows and drifted
evaluation windows should be calibrated separately rather than pooled into a
single threshold rule. This is exactly the discipline used in the logged
benchmark protocol, where thresholds are fit on healthy windows only and then
evaluated on later regimes. The same line of work also supports
bootstrap-based uncertainty quantification for Sinkhorn-type functionals, which
matches the paired bootstrap intervals reported for KuaiRand.
\end{remark}

\begin{assumption}[$W_2$-Lipschitz Drift]\label{ass:adrift}
The oracle trajectory $t \mapsto \Pot^*(S,t)$ is $W_2$-Lipschitz with drift
rate $\zeta \geq 0$: for all $s,t \geq 0$,
\[
  W_2\!\left(\Pot^*(S,s),\,\Pot^*(S,t)\right)
  \;\leq\; \zeta\,|t-s|.
\]
\end{assumption}

For the Gaussian family with common covariance, \cref{ass:adrift} reduces to a
Lipschitz bound on the mean trajectory.

\begin{proposition}[Estimation Error Bound for $\hat\sigma_P$]%
\label{prop:tci_estimation_error}
Let \cref{ass:adrift} hold, and assume each oracle marginal $\Pot^*(S,t)$ is
$K$-subgaussian. Let $\hat\Pot^*$ be the sliding-window empirical estimator
with window $n$ (\cref{def:potestimator}), and let
$\tau = (n-1)/2$ be the effective lag of the centered window.
Then under the entropic regularization of \cref{rem:sinkhorn},
\[
  \mathbb{E}\,\bigl|\sigma_P(t) - \hat\sigma_P(t)\bigr|
  \;\leq\;
  C_K\,n^{-1/2} + \zeta\tau,
\]
where $C_K$ absorbs the finite-sample Sinkhorn constant, the fixed
regularization $\varepsilon$, the ambient geometry, and the uniform subgaussian
envelope. This is a finite-memory upper bound for the sliding-window class, not
a minimax statement over all adaptive estimators.
\end{proposition}

\begin{proof}
Let
\[
  \bar\Pot^*(S,t)
  \;\coloneqq\;
  \frac{1}{n}\sum_{k=t-n+1}^{t}\Pot^*(S,k)
\]
be the oracle window-average law. The finite-sample term is the usual
subgaussian Sinkhorn estimation error: by \citet[Thm.~3]{GenevayEtAl2019}, the
debiased divergence between $\bar\Pot^*(S,t)$ and its empirical estimator is of
order $n^{-1/2}$, with the regularization dependence absorbed into $C_K$.

The additional error comes from temporal staleness of the window. By convexity
of $W_2^2$ in its second argument,
\[
  W_2\!\left(\Pot^*(S,t),\bar\Pot^*(S,t)\right)
  \leq
  \frac{1}{n}\sum_{k=t-n+1}^{t}
  W_2\!\left(\Pot^*(S,t),\Pot^*(S,k)\right).
\]
Applying \cref{ass:adrift},
\begin{align*}
  W_2\!\left(\Pot^*(S,t),\bar\Pot^*(S,t)\right)
  &\leq \frac{\zeta}{n}\sum_{k=t-n+1}^{t}|t-k| \\
  &= \zeta\,\frac{n-1}{2}
   = \zeta\tau.
\end{align*}
This is the explicit lag-bias term. Combining it with the
$C_K n^{-1/2}$ statistical term yields the stated decomposition. Finally, the
score map $x \mapsto 1/(1+\tfrac{1}{2}x^2)$ is globally Lipschitz on
$\mathbb{R}_{\geq 0}$, so the same order transfers to
$\mathbb{E}|\sigma_P(t)-\hat\sigma_P(t)|$.
\end{proof}

\begin{remark}[Tightness of the lag-bias term]\label{rem:lag_tight}
The lag term is not merely an artifact of upper-bounding. For monotone
translation trajectories $\Pot^*(S,t)=\Pot_0(S-vt)$ with $\|v\|=\zeta$,
the current law and the centered window-average law are separated by exactly the
mean age of the window, so the temporal bias is of order $\zeta\tau$.
The cube-root balance therefore reflects a real finite-memory limitation rather
than a proof-only looseness.
\end{remark}

\Cref{prop:tci_estimation_error} grounds $\ScoreHat$: the estimation error
decreases with window size and increases with drift. This same
$n^{-1/2}$ versus $\zeta\tau$ trade-off yields the tracking limit.
When window lag is accounted for explicitly ($\tau = n/2$ for a centred
window), the total estimation error decomposes as a sum of two opposing terms.
Minimising over $n$ yields the cube-root optimum for the averaging class,
including the sliding-window and exponentially weighted schemes.

\begin{definition}[Lag-Inclusive Estimation Error]\label{def:lag_error}
For a sliding window of size $n$ centred at lag $\tau = n/2$, the
total expected estimation error is:
\[
  \mathcal{E}(n) \;=\; C_K\, n^{-1/2} \;+\; \tfrac{1}{2}\,\zeta\, n,
\]
where the first term is the \emph{variance} (decreasing in $n$) and the
second is the \emph{lag-bias} (increasing in $n$).
\end{definition}

\begin{proposition}[Optimal Window Under Drift]%
\label{prop:coherence_delay}
Under \cref{def:lag_error}, the estimation error $\mathcal{E}(n)$ is
minimised at the \emph{optimal window}
\[
  n^* \;=\; \left(\frac{C_K}{\zeta}\right)^{2/3},
\]
yielding the \emph{minimum achievable tracking error}
\[
  \mathcal{E}_{\min} \;=\; \tfrac{3}{2}\,C_K^{2/3}\,\zeta^{1/3}.
\]
\end{proposition}

\begin{proof}
Setting $d\mathcal{E}/dn = 0$:
$-\tfrac{1}{2}C_K n^{-3/2} + \tfrac{1}{2}\zeta = 0$,
giving $n^* = (C_K/\zeta)^{2/3}$.
Substituting back:
$\mathcal{E}(n^*) = C_K(C_K/\zeta)^{-1/3} + \tfrac{1}{2}\zeta(C_K/\zeta)^{2/3}
= C_K^{2/3}\zeta^{1/3} + \tfrac{1}{2}C_K^{2/3}\zeta^{1/3}
 = \tfrac{3}{2}C_K^{2/3}\zeta^{1/3}$.
\end{proof}

\begin{corollary}[EWMA Analogue]\label{cor:ewma_cube_root}
Let the finite-memory estimator use exponential weights
$w_j = \alpha(1-\alpha)^j$ for $j \geq 0$, so its effective lag is
$\tau_\alpha = \sum_{j\geq 0} j w_j = (1-\alpha)/\alpha$ and its effective
sample size is
$n_{\mathrm{eff}}(\alpha) = 1 / \sum_{j\geq 0} w_j^2 = (2-\alpha)/\alpha$.
Under \cref{ass:adrift}, the resulting finite-memory error obeys
\[
  \mathcal{E}_{\mathrm{EWMA}}(\alpha)
  \;\leq\;
  C_K\sqrt{\frac{\alpha}{2-\alpha}}
  + \zeta\,\frac{1-\alpha}{\alpha}.
\]
Hence, in the small-$\alpha$ regime,
$\alpha^* \propto (\zeta/C_K)^{2/3}$ and
$\mathcal{E}_{\mathrm{EWMA},\min} \propto C_K^{2/3}\zeta^{1/3}$.
\end{corollary}

\begin{proof}
The variance term is the sliding-window rate with $n$ replaced by the effective
sample size $n_{\mathrm{eff}}(\alpha)$. The lag term is the
$W_2$-Lipschitz drift rate multiplied by the effective age
$\tau_\alpha$. For small $\alpha$ this becomes
$\mathcal{E}_{\mathrm{EWMA}}(\alpha)
= O(C_K\alpha^{1/2} + \zeta\alpha^{-1})$.
Balancing the two exponents gives
$\alpha^* = O((\zeta/C_K)^{2/3})$, and substitution yields the same
$\zeta^{1/3}$ minimum-error law as the sliding-window case.
\end{proof}

\begin{definition}[Causal Averaging Kernel]\label{def:avg_kernel}
A \emph{causal averaging kernel} is a sequence of weights
$w = (w_0,w_1,\dots)$ such that $w_j \geq 0$ and $\sum_{j\geq 0} w_j = 1$.
Its effective size and effective lag are
\[
  n_{\mathrm{eff}}(w) \coloneqq \frac{1}{\sum_{j\geq 0} w_j^2},
  \qquad
  \tau_{\mathrm{eff}}(w) \coloneqq \sum_{j\geq 0} j w_j.
\]
\end{definition}

\begin{proposition}[Averaging-Kernel Lower Bound]\label{prop:kernel_lower_bound}
Let $X_t = \mu_t + \xi_t$ with $\xi_t \sim \mathcal{N}(0,\sigma^2)$ i.i.d.
and let the drift trajectory be the monotone translation model
$\mu_t = \zeta t$ with $\zeta \geq 0$. For any causal averaging kernel
$w$ as in \cref{def:avg_kernel}, the kernel estimator
\[
  \hat\mu_t^{(w)} \coloneqq \sum_{j\geq 0} w_j X_{t-j}
\]
satisfies
\[
  \sqrt{\mathbb{E}\,\bigl(\hat\mu_t^{(w)} - \mu_t\bigr)^2}
  \;\geq\;
  c_1\,\sigma\,n_{\mathrm{eff}}(w)^{-1/2}
  \vee
  c_2\,\zeta\,\tau_{\mathrm{eff}}(w)
\]
for universal constants $c_1,c_2>0$.
Equivalently, the mean-squared error obeys the exact decomposition
\[
  \mathbb{E}\,\bigl(\hat\mu_t^{(w)} - \mu_t\bigr)^2
  \;=\;
  \sigma^2 n_{\mathrm{eff}}(w)^{-1} + \zeta^2 \tau_{\mathrm{eff}}(w)^2,
\]
and therefore the RMSE lower bound. In particular, optimizing over the averaging
family yields the same cube-root exponent as \cref{prop:coherence_delay}.
\end{proposition}

\begin{proof}
The estimator error decomposes as
\[
  \hat\mu_t^{(w)} - \mu_t
  = \underbrace{\zeta\sum_{j\geq 0} j w_j}_{\text{bias}}
    + \underbrace{\sum_{j\geq 0} w_j \xi_{t-j}}_{\text{noise}}.
\]
Hence the bias magnitude is exactly $\zeta\tau_{\mathrm{eff}}(w)$ and the noise
term is Gaussian with variance $\sigma^2\sum_{j\geq 0}w_j^2 = \sigma^2 /
n_{\mathrm{eff}}(w)$. Therefore
\[
  \mathbb{E}\,\bigl(\hat\mu_t^{(w)} - \mu_t\bigr)^2
  = \zeta^2 \tau_{\mathrm{eff}}(w)^2 + \sigma^2 n_{\mathrm{eff}}(w)^{-1},
\]
and taking square roots gives the stated lower bound, since
\(\sqrt{a+b} \geq a^{1/2} \vee b^{1/2}\) for nonnegative $a,b$.
Optimizing the resulting trade-off over the averaging family recovers the same
$2/3$ window scaling and $1/3$ error exponent as the upper bound.
\end{proof}

The kernel result is exact but still tied to linear averaging. The minimax
problem removes that linearity assumption on the estimator side by moving to all
measurable rules that only observe a finite recent window. The price is that the
signal class is now the full discrete
$\zeta$-Lipschitz family rather than the single monotone translation witness.

\begin{proposition}[Window-Restricted Minimax Lower Bound]
\label{prop:window_minimax}
Fix an integer $m \geq 1$ and suppose the estimator only observes the last $m$
samples
\[
  X_{t-m+1},\dots,X_t,
  \qquad
  X_s = \mu_s + \xi_s,
  \qquad
  \xi_s \stackrel{\mathrm{i.i.d.}}{\sim} \mathcal{N}(0,\sigma^2).
\]
Let $\mathcal{L}_m(\zeta)$ be the discrete $\zeta$-Lipschitz class
\[
  \mathcal{L}_m(\zeta)
  \coloneqq
  \left\{(\mu_{t-m+1},\dots,\mu_t):
    |\mu_i-\mu_j| \leq \zeta |i-j|\ \text{for all}\ i,j
  \right\}.
\]
For any measurable estimator $\hat\mu_t = \delta_m(X_{t-m+1:t})$, define the
worst-case root risk
\[
  R_m(\delta_m)
  \coloneqq
  \sup_{\mu \in \mathcal{L}_m(\zeta)}
  \Bigl(\mathbb{E}_\mu (\hat\mu_t - \mu_t)^2\Bigr)^{1/2}.
\]
Then there exists a universal constant $c>0$ such that
\[
  \inf_{\delta_m} R_m(\delta_m)
  \;\geq\;
  c\,\max\!\left\{
    \sigma m^{-1/2},\;
    \sup_{1\leq h\leq m}\min\!\left(\sigma h^{-1/2},\; \zeta h\right)
  \right\}.
\]
In particular, if $1 \leq (\sigma/\zeta)^{2/3} \leq m$, then
\[
  \inf_{\delta_m} R_m(\delta_m)
  \;\geq\;
  c\,\sigma^{2/3}\zeta^{1/3}.
\]
\end{proposition}

\begin{proof}
Write $R_m^* \coloneqq \inf_{\delta_m} R_m(\delta_m)$.

We exhibit two independent obstructions. The first is the ordinary
finite-sample estimation floor, obtained by restricting to the constant-path
subclass $\mu_s \equiv \theta$, which is contained in
$\mathcal{L}_m(\zeta)$. For the two hypotheses
$\theta_\pm = \pm \sigma /(2\sqrt{m})$, the induced product-Gaussian laws have
Kullback--Leibler divergence
\[
  \mathrm{KL}(P_+,P_-)
  = \frac{1}{2\sigma^2} \sum_{j=1}^{m} (\theta_+ - \theta_-)^2
  = \frac{1}{2\sigma^2} m\left(\frac{\sigma}{\sqrt{m}}\right)^2
  = \frac{1}{2}.
\]
For any estimator $\delta_m$, the plug-in test $\psi = \mathbf{1}\{\delta_m \geq 0\}$
must err whenever the estimate falls on the wrong side of the origin; on that
event the endpoint error is at least $\sigma /(2\sqrt{m})$. By Le Cam's
two-point lemma and Pinsker's inequality, the average testing error is at least
\[
  \frac{1-\mathrm{TV}(P_+,P_-)}{2}
  \geq
  \frac{1-\sqrt{\mathrm{KL}(P_+,P_-)/2}}{2}
  \geq \frac{1}{4}.
\]
Hence
\[
  R_m^* \;\geq\; \frac{\sigma}{8\sqrt{m}}
\]
so one valid explicit choice is $c_0 = 1/8$.

The second obstruction is local drift. Fix any $1\leq h\leq m$ and define
\[
  \beta_h \coloneqq \min\!\left\{\zeta,\; \frac{\sigma}{4h^{3/2}}\right\},
  \qquad
  \mu^{\pm,h}_{t-j} \coloneqq \pm \,\beta_h\,(h-j)_+,
  \qquad 0\leq j\leq m-1,
\]
where $(u)_+ \coloneqq \max\{u,0\}$. Each path belongs to
$\mathcal{L}_m(\zeta)$ because its slope magnitude is at most $\beta_h \leq
\zeta$. The endpoint gap is
\[
  |\mu_t^{+,h} - \mu_t^{-,h}| = 2\beta_h h.
\]
The corresponding Gaussian product laws satisfy
\[
  \mathrm{KL}(P_{+,h},P_{-,h})
  = \frac{1}{2\sigma^2} \sum_{j=0}^{m-1}
    \bigl(2\beta_h (h-j)_+\bigr)^2
  \leq \frac{2\beta_h^2}{\sigma^2} \sum_{r=1}^{h} r^2
  \leq \frac{2\beta_h^2}{\sigma^2} h^3
  \leq \frac{1}{8}.
\]
Applying the same reduction from estimation to testing, now with endpoint error
threshold $\beta_h h$, and using Le Cam together with Pinsker gives average
testing error at least
\[
  \frac{1-\mathrm{TV}(P_{+,h},P_{-,h})}{2}
  \geq
  \frac{1-\sqrt{\mathrm{KL}(P_{+,h},P_{-,h})/2}}{2}
  \geq \frac{3}{8}.
\]
Therefore
\[
  R_m^* \;\geq\; \frac{3}{8}\,\beta_h h
  \;\geq\;
  \frac{3}{32}\,\min\!\left\{\zeta h,\; \sigma h^{-1/2}\right\}
\]
which is again a fully explicit numerical lower bound. Since $h$ was arbitrary,
\[
  R_m^* \;\geq\;
  \frac{3}{32}\,\sup_{1\leq h\leq m}
  \min\!\left\{\zeta h,\; \sigma h^{-1/2}\right\}.
\]

Taking the maximum of the constant-path and ramp bounds proves the displayed
lower bound. For the cube-root regime, assume
$1 \leq h_\star \coloneqq (\sigma/\zeta)^{2/3} \leq m$. Choose an integer
$\bar h \in [h_\star/2, h_\star]$. Then
\[
  \zeta \bar h \geq \frac{1}{2}\,\zeta h_\star
  = \frac{1}{2}\,\sigma^{2/3}\zeta^{1/3},
\]
and
\[
  \sigma \bar h^{-1/2}
  \geq \sigma h_\star^{-1/2}
  = \sigma^{2/3}\zeta^{1/3}.
\]
Therefore
\[
  \min\!\left\{\sigma \bar h^{-1/2},\; \zeta \bar h\right\}
  \geq \frac{1}{2}\,\sigma^{2/3}\zeta^{1/3},
\]
and the ramp bound gives
\[
  R_m^* \geq \frac{3}{64}\,\sigma^{2/3}\zeta^{1/3}
\]
so one admissible explicit constant in the critical-window regime is
$c = 3/64$. This proves the claim.
\end{proof}

In other words, the displayed cube-root floor is operational exactly in the
critical-window regime $1 \leq (\sigma/\zeta)^{2/3} \leq m$: outside that range,
the problem is either noise-dominated or too window-limited for the minimax
obstruction to reach the balancing scale.

\begin{remark}[Interpretation: a finite-memory floor]\label{rem:cdfloor}
\Cref{prop:coherence_delay} says that continuous drift imposes a nonzero floor
for the finite-memory averaging class: $\mathcal{E}_{\min} \propto \zeta^{1/3}$
and $n^* \propto \zeta^{-2/3}$. In practice this becomes a self-tuning rule,
$n_t^* = (C_K/\hat\zeta_t)^{2/3}$, with the critical-window lower bound in
\cref{prop:window_minimax} showing that the same obstruction persists for any
estimator restricted to the last $m$ samples.
\end{remark}
